\subsection{Purpose of Mathematical sets}

Sets are among the most fundamental objects in mathematics. They provide a mechanism for mathematicians and engineers to concisely analyze, reference, and categorize various objects. In fact, a set is essentially an abstract container.

A set is effectively the primary encapsulation tool for mathematics. Imagine creating a program that requires a function with many parameters. Functions with a large number of parameters become difficult to manage. Programmers often wrap the parameters in a struct or class. As a result, the function can accept one parameter (the struct or class) and let the function unpack the contents. Mathematicians similarly encapsulate math objects in a set, give the set a name, and then reference the set using its name or identifier.

The example below represents a very traditional way of representing sets. Notice the curly braces are the common notation when creating a set.

\[
 A = \{4, 6, 1, -3, 14, 99\}
\]

In the example above, a set of numbers has been created and labeled $A$. Once declared, the letter/symbol $A$ can be used to refer to this collection of values. The elements of a set can be essentially anything, including other sets. Consider the set $B$ composed of colors.

\[
 B = \{\text{blue},\; \text{red},\; \text{black}\}
\]

Set $C$ (below) is a set where the elements are also sets.

\[
 C = \{\; \{1,2,3,4,5\},\; \{\pi, 14, 2\}\;\}
\]

